\chapter*{Заключение}
\addcontentsline{toc}{chapter}{Заключение}

В настоящей работе разработан и исследован алгоритм приоритизации тестовых случаев GA-FP+Div, 
представляющий собой генетический алгоритм с улучшенной многокомпонентной фитнесс-функцией. 
Фитнесс-функция учитывает три нормализованных критерия: покрытие кода, вероятность дефектов на основе 
CIA-анализа и структурное разнообразие тестов по расстоянию Жаккара.

В ходе работы получены следующие результаты:

\begin{enumerate}
    \item Проведён анализ существующих методов приоритизации тестовых случаев, выявлены ограничения жадных 
    и эволюционных подходов, обоснована необходимость многокритериальной фитнесс-функции.

    \item Разработана фитнесс-функция, объединяющая покрытие кода, вероятность дефектов и структурное 
    разнообразие тестов. Реализованы операторы OX-кроссовера, адаптивной мутации и элитизма.

    \item Выполнена программная реализация на языке Java с поддержкой загрузки данных в форматах 
    BigFaultMatrix и Defects4J, а также сохранением результатов в форматах CSV и JSON.

    \item Проведён сравнительный эксперимент на наборе данных BigFaultMatrix (4 вариации) и 7 проектах 
    Defects4J. Алгоритм GA-FP+Div показал наилучшие результаты на 5 из 7 проектов Defects4J, 
    превзойдя базовый генетический алгоритм на всех проектах.

    \item Установлено что эффективность GA-FP+Div определяется качеством CIA-данных: на проектах с 
    богатым CIA-анализом (Codec/18, Lang/65) преимущество алгоритма наиболее выражено.
\end{enumerate}

Направления дальнейших исследований включают: автоматический подбор весов фитнесс-функции методами 
настройки гиперпараметров; расширение бенчмарка за счёт дополнительных проектов Defects4J; исследование 
влияния параметра $k$ на качество приоритизации.